\chapter{Как делать свою прошивку}\label{ch:svoya}

Цель всей работы --- понять, что можно менять. Мы теперь знаем достаточно, чтобы описать три уровня вмешательства от безопасного к рискованному.

\section{Три уровня}

\begin{figure}[tbp]
\centering
\begin{tikzpicture}[
  lv/.style={draw=rule, fill=paper, rounded corners=3pt, minimum width=118mm, minimum height=15mm, text width=112mm, align=left, font=\sffamily\small, inner sep=5pt},
  l1/.style={lv, draw=ledgreen, fill=ledgreen!8},
  l2/.style={lv, draw=ledyellow, fill=ledyellow!10},
  l3/.style={lv, draw=ledred, fill=ledred!6}
]
\node[l1] (a) {\textbf{Уровень 1. Данные вместо кода.} Менять рисунок ковра, другое семейство меток в программе модуля, параметры автопилота. Дрон трогать не нужно.};
\node[l2, below=4mm of a] (b) {\textbf{Уровень 2. Своя камера вместо модуля.} Любая плата, которая шлёт автопилоту кадр из главы~\ref{ch:uart}. Позволяет ArUco, QR и всё остальное.};
\node[l3, below=4mm of b] (c) {\textbf{Уровень 3. Свой автопилот.} Замена \texttt{pilot\_f}. Огромный проект: основой служат карта (\texttt{melis}) и логика (\texttt{pilot}).};
\end{tikzpicture}
\caption{Лестница вмешательства: чем ниже, тем больше риск.}\label{fig:ladder}
\end{figure}

\section{Уровень 1: изменить программу модуля}

Самый короткий путь к другому семейству меток. Что известно: маска семейств лежит в функции \texttt{imlib\_find\_apriltag\_app\_load} как константа \texttt{0x10}; остальные четыре семейства скомпилированы. Обновление идёт через USB-диск модуля. \pr{}

\begin{enumerate}
  \item Снять полную копию модуля (она у нас есть).
  \item Изменить константу в копии \texttt{libmaix\_ckh831.so}.
  \item Положить папку \texttt{maix\_dist} на USB-диск и перезапустить модуль.
  \item Проверить на дроне без пропеллеров, что дрон видит номер метки.
\end{enumerate}

\begin{vnimanie}[Ограничения]
Размер метки и калибровка камеры в программе модуля пока не найдены. Для другого семейства может понадобиться другой размер метки. Если программа падает четыре раза подряд, модуль сбросится.
\end{vnimanie}

\section{Уровень 2: своя камера}

Автопилот ничего не знает про меток: ему нужен только кадр. Свою камеру можно сделать из любой платы с UART, если она умеет:

\begin{itemize}
  \item слать кадр из главы~\ref{ch:uart} с \textbf{правильным заголовком для версии автопилота} (один байт \texttt{CD} для версии марта 2023, два для 1.2.8);
  \item отвечать на команду проверки связи \texttt{9} (ответ \texttt{F9 66});
  \item понимать команду выбора режима \texttt{1}, хотя бы режим 0;
  \item при желании управлять светодиодами и ШИМ.
\end{itemize}

Дальномер ToF (байты 18--19 кадра) нужно либо тоже поставить, либо слать нули; как автопилот использует это поле, мы не выяснили. \nei{}

\begin{zametka}[Что даёт ArUco]
Алгоритм ArUco из OpenCV находит метки другого семейства и даёт те же номера, смещения и угол. Если свой модуль посчитает их в единицах, которые ждёт автопилот (раздел~\ref{sec:units}), дрон полетит, как и по штатным меткам. \vyv{}
\end{zametka}

\section{Единицы}\label{sec:units}

Для расчёта смещений в единицах автопилота (по формуле главы~\ref{ch:uart} и таблице ниже):

\begin{tabularx}{\textwidth}{@{}l Y@{}}
\toprule
\hdr{Поле} & \hdr{Единица (приблизительно)} \\ \midrule
$x$, $y$ & 0,83 мм на единицу ($8{,}333\cdot 10^{-4}$ м). \\
Расстояние & 1,02 мм на единицу ($10{,}2\cdot 10^{-4}$ м). \\
Угол & 0,1\textdegree. \\
Номер & Целое от 0 до 287. \\ \bottomrule
\end{tabularx}

\vspace{4pt}
\vyv{} Единицы выведены из формулы автопилота и подтверждены одним живым кадром; для точности нужна серия замеров на разных высотах.

\section{Уровень 3: свой автопилот}

Полная замена автопилота --- отдельный большой проект. Что у нас есть для него: формат всех кадров, структуры данных, режимы и их порядок, раскладка миссии, константы формулы метки, карта прошивки C10. Чего нет: исходного кода регуляторов, драйверов датчиков, закрытого алгоритма оптического потока, загрузчика и схем платы. \nei{}

\section{Чек-лист перед первым полётом}

\begin{enumerate}
  \item Полная копия флеша дрона и программы модуля.
  \item Проверка без пропеллеров: приём кадра (счётчик \texttt{rx}, буфер автопилота).
  \item Проверка на привязи, на малой высоте.
  \item Только потом свободный полёт и только на запасном дроне.
\end{enumerate}

\begin{opasno}
Ошибка в данных камеры напрямую влияет на управление полётом: дрон поверит неверному положению. Любое изменение в цепочке метка $\to$ кадр $\to$ автопилот проверяйте на столе, без пропеллеров.
\end{opasno}
